% LaTeX Template for AI Problem Analysis Output
% 2017 AMC 12A Problem 17 Strategic Analysis
% IMPORTANT: Use LaTeX commands for all mathematical symbols, NOT unicode characters

\documentclass[]{article}
\usepackage[utf8]{inputenc}
\usepackage{amsmath, amssymb, amsthm}
\usepackage{geometry}
\usepackage{xcolor}
\usepackage[most]{tcolorbox}
\usepackage{enumitem}
\usepackage{fancyhdr}

% Page setup
\geometry{margin=1in}
\pagestyle{fancy}
\fancyhf{}
\rhead{AMC12A 2017 Problem 17 Analysis}
\lhead{\today}
\cfoot{\thepage}

% Color definitions
\definecolor{problemconfig}{RGB}{230, 242, 255}    % Light blue
\definecolor{strategic}{RGB}{255, 230, 242}       % Light pink
\definecolor{tactical}{RGB}{242, 255, 230}        % Light green
\definecolor{techniques}{RGB}{255, 242, 230}      % Light yellow
\definecolor{verification}{RGB}{255, 230, 230}    % Light red
\definecolor{solution}{RGB}{230, 255, 255}        % Light cyan
\definecolor{competition}{RGB}{242, 242, 255}     % Light purple
\definecolor{gold}{RGB}{218, 165, 32}

% Box styles
\newtcolorbox{problemconfigbox}{
    colback=problemconfig,
    colframe=blue,
    title=Problem Configuration Analysis,
    fonttitle=\bfseries,
    arc=3pt,
    boxrule=1pt,
    breakable
}

\newtcolorbox{strategicbox}{
    colback=strategic,
    colframe=purple,
    title=Strategic Framework Application,
    fonttitle=\bfseries,
    arc=3pt,
    boxrule=1pt,
    breakable
}

\newtcolorbox{tacticalbox}{
    colback=tactical,
    colframe=green,
    title=Tactical Methodology Selection,
    fonttitle=\bfseries,
    arc=3pt,
    boxrule=1pt,
    breakable
}

\newtcolorbox{techniquesbox}{
    colback=techniques,
    colframe=orange,
    title=Conversion Techniques Application,
    fonttitle=\bfseries,
    arc=3pt,
    boxrule=1pt,
    breakable
}

\newtcolorbox{verificationbox}{
    colback=verification,
    colframe=red,
    title=Verification Strategy Design,
    fonttitle=\bfseries,
    arc=3pt,
    boxrule=1pt,
    breakable
}

\newtcolorbox{solutionbox}{
    colback=solution,
    colframe=cyan,
    title=Solution Roadmap,
    fonttitle=\bfseries,
    arc=3pt,
    boxrule=1pt,
    breakable
}

\newtcolorbox{competitionbox}{
    colback=competition,
    colframe=purple,
    title=Competition Strategy Integration,
    fonttitle=\bfseries,
    arc=3pt,
    boxrule=1pt,
    breakable
}

\newtcolorbox{keyinsightbox}{
    colback=yellow!20,
    colframe=gold,
    title=Key Insight,
    fonttitle=\bfseries,
    arc=3pt,
    boxrule=2pt,
    leftrule=5pt,
    breakable
}

\newtcolorbox{warningbox}{
    colback=red!10,
    colframe=red!50,
    title=Warning: Common Pitfall,
    fonttitle=\bfseries,
    arc=3pt,
    boxrule=2pt,
    leftrule=5pt,
    breakable
}

\newtcolorbox{tipbox}{
    colback=green!10,
    colframe=green!50,
    title=Pro Tip,
    fonttitle=\bfseries,
    arc=3pt,
    boxrule=2pt,
    leftrule=5pt,
    breakable
}

\begin{document}

\title{AMC12A 2017 Problem 17: Strategic Analysis}
\author{Complex Numbers \& Roots of Unity}
\date{\today}
\maketitle

\section*{Problem Statement}
There are $24$ different complex numbers $z$ such that $z^{24}=1$. For how many of these is $z^6$ a real number?

\textbf{Answer Choices:} $\textbf{(A)}\ 0 \qquad\textbf{(B)}\ 4 \qquad\textbf{(C)}\ 6 \qquad\textbf{(D)}\ 12 \qquad\textbf{(E)}\ 24$

\begin{problemconfigbox}
\subsection*{A. Problem Classification}
\begin{itemize}[leftmargin=*]
    \item \textbf{Problem Type}: To Find - counting the number of 24th roots of unity with specific property
    \item \textbf{Difficulty Band}: Late-band (Problem 17 of 25) - requires knowledge of complex numbers and roots of unity
    \item \textbf{Competition Context}: AMC12 (75 min, 25 problems) - approximately 3-4 minutes recommended for this problem
    \item \textbf{Mathematical Domain}: Complex Numbers, Algebra (roots of unity, exponential form, De Moivre's Theorem)
\end{itemize}

\subsection*{B. Problem Structure Identification}
\begin{itemize}[leftmargin=*]
    \item \textbf{Given Information}: 
    \begin{itemize}
        \item There are 24 complex numbers satisfying $z^{24} = 1$ (the 24th roots of unity)
        \item We need these 24 values of $z$
        \item For each $z$, we compute $z^6$
    \end{itemize}
    \item \textbf{Constraints}: 
    \begin{itemize}
        \item $z$ must be a 24th root of unity
        \item $z^6$ must be a real number (imaginary part equals 0)
    \end{itemize}
    \item \textbf{Objective}: Count how many of the 24 roots of unity have $z^6$ real
    \item \textbf{Hidden Assumptions}: 
    \begin{itemize}
        \item Knowledge of roots of unity: $z = e^{2\pi i k/24}$ for $k = 0, 1, 2, \ldots, 23$
        \item Real numbers in complex form have argument $0$ or $\pi$ (multiples of $\pi$)
        \item Connection to De Moivre's Theorem: $z^n = e^{in\theta}$
    \end{itemize}
    \item \textbf{Answer Choice Leverage}: 
    \begin{itemize}
        \item Answer choices range from 0 to 24, with intermediate values 4, 6, 12
        \item Choice (D) 12 is exactly half of 24, suggesting parity/symmetry argument
        \item Choice (C) 6 is exactly one-fourth, choice (B) 4 is one-sixth
        \item The even distribution suggests a divisibility or modular arithmetic pattern
    \end{itemize}
\end{itemize}

\subsection*{C. Strategic Orientation}
\begin{itemize}[leftmargin=*]
    \item \textbf{Hypothesis}: The 24th roots of unity are $z_k = e^{2\pi i k/24} = e^{i\pi k/12}$ for $k = 0, 1, \ldots, 23$
    \item \textbf{Conclusion}: Determine which values of $k$ make $z_k^6$ real
    \item \textbf{Notation}: 
    \begin{itemize}
        \item $z_k = e^{i\pi k/12}$ for the $k$-th root of unity
        \item $z_k^6 = e^{i\pi k/2}$ using exponent rules
        \item A complex number $e^{i\theta}$ is real when $\theta$ is a multiple of $\pi$
    \end{itemize}
    \item \textbf{Intuitive Approach}: 
    \begin{itemize}
        \item The 24th roots form a regular 24-gon on the unit circle
        \item Raising to the 6th power rotates each point by 6 times its angle
        \item Real numbers on the unit circle are at $\pm 1$ (angles $0$ and $\pi$)
    \end{itemize}
    \item \textbf{Cross-Domain Potential}: 
    \begin{itemize}
        \item \textbf{Algebraic approach}: Direct computation with $z^6 = e^{i\pi k/2}$, checking when imaginary part vanishes
        \item \textbf{Number theory approach}: Modular arithmetic - when is $k/2$ an integer? When $k$ is even
        \item \textbf{Geometric approach}: Visual representation on the unit circle, identifying which roots map to real axis
        \item \textbf{Group theory approach}: $z^6$ being real means $z^6 \in \{1, -1\}$, which are the 2nd roots of unity
    \end{itemize}
\end{itemize}
\end{problemconfigbox}

\begin{strategicbox}
\subsection*{A. Psychological Strategies}
\begin{itemize}[leftmargin=*]
    \item \textbf{Mental Toughness}: This is a late-band problem requiring confidence with complex numbers. If roots of unity are unfamiliar, acknowledge this and either (1) attempt simpler approach via answer choices, or (2) skip and return if time permits
    \item \textbf{Creativity Cultivation}: Consider multiple representations - exponential form $e^{i\theta}$, trigonometric form $\cos\theta + i\sin\theta$, or geometric interpretation on the unit circle
    \item \textbf{Iterative Refinement}: Start with small cases - what if the problem asked about $z^{12} = 1$? Or if we needed $z^3$ real? Pattern recognition from simpler versions
\end{itemize}

\subsection*{B. Investigation Strategies}
\begin{itemize}[leftmargin=*]
    \item \textbf{Startup Orientation}: Recognize this as a roots of unity problem. The 24 solutions to $z^{24} = 1$ are evenly distributed on the unit circle
    \item \textbf{Get Your Hands Dirty}: Write out the first few roots explicitly:
    \begin{itemize}
        \item $k=0$: $z = 1$, so $z^6 = 1$ (real) \checkmark
        \item $k=1$: $z = e^{i\pi/12}$, so $z^6 = e^{i\pi/2} = i$ (not real)
        \item $k=2$: $z = e^{i\pi/6}$, so $z^6 = e^{i\pi} = -1$ (real) \checkmark
        \item $k=3$: $z = e^{i\pi/4}$, so $z^6 = e^{3i\pi/2} = -i$ (not real)
        \item $k=4$: $z = e^{i\pi/3}$, so $z^6 = e^{2i\pi} = 1$ (real) \checkmark
    \end{itemize}
    Pattern emerges: real when $k$ is even!
    \item \textbf{Penultimate Step}: The final step is counting even values of $k$ in $\{0,1,2,\ldots,23\}$
    \item \textbf{Wishful Thinking}: "I wish I could just identify when $e^{i\pi k/2}$ is real" - this happens exactly when $\pi k/2$ is a multiple of $\pi$, i.e., when $k$ is even
    \item \textbf{Make It Easier}: Replace $z^{24}=1$ with $z^{12}=1$ or $z^6=1$ to see the pattern more clearly. For $z^{12}=1$, when is $z^6$ real? This gives insight into the structure
    \item \textbf{Conjecture via Small Cases}: As computed above, $k$ even gives real values. Test this: there are 12 even values in $\{0,1,\ldots,23\}$
    \item \textbf{Visualization}: Draw the 24th roots as points on the unit circle. The condition $z^6$ real means rotating by angle $6 \cdot 2\pi k/24 = \pi k/2$ lands on the real axis. This happens at angles $0, \pi, 2\pi, 3\pi, \ldots$ which correspond to $k = 0, 2, 4, 6, \ldots$
\end{itemize}

\subsection*{C. Late-Band Strategic Playbook (Problems 17-25)}
\begin{itemize}[leftmargin=*]
    \item \textbf{Answer-Choice Leverage}: Answer choice (D) is 12, exactly half of 24. This strongly suggests an even/odd parity argument. If every other root works, we get exactly half
    \item \textbf{Make It Easier}: Consider $z^8 = 1$ instead. Then $z = e^{2\pi i k/8}$ for $k=0,\ldots,7$, and $z^4 = e^{i\pi k}$. This is real when $k$ is even, giving 4 out of 8 values. Pattern: half!
    \item \textbf{Penultimate Step}: Working backwards - if the answer is 12, what pattern would give 12? Even values of $k$: there are exactly 12 even numbers in $\{0,1,\ldots,23\}$
    \item \textbf{Wishful Thinking}: Assume $z^6 \in \{1, -1\}$ (the only real numbers on the unit circle). Then $z^{12} = (z^6)^2 = 1$, so $z$ is a 12th root of unity. But we need all 24th roots where this holds, which means every other 24th root
    \item \textbf{Conjecture via Small Cases}: From experimental computation, pattern is clear: $k$ even $\Leftrightarrow$ $z^6$ real
\end{itemize}

\subsection*{D. Argument Construction Strategy}
\begin{itemize}[leftmargin=*]
    \item \textbf{Direct Proof}: 
    \begin{itemize}
        \item The 24th roots of unity are $z = e^{2\pi i k/24}$ for $k = 0, 1, \ldots, 23$
        \item Then $z^6 = e^{2\pi i k \cdot 6/24} = e^{i\pi k/2}$
        \item For $z^6$ to be real, we need $\sin(\pi k/2) = 0$, which occurs when $\pi k/2 = n\pi$ for some integer $n$
        \item This gives $k = 2n$, i.e., $k$ must be even
        \item There are 12 even values in $\{0, 1, 2, \ldots, 23\}$
    \end{itemize}
    \item \textbf{Proof by Contradiction}: Not applicable here (direct proof is straightforward)
    \item \textbf{Mathematical Induction}: Not applicable (this is a counting problem, not a sequence)
    \item \textbf{Stylistic Conventions}: WLOG not needed; the argument is direct and complete
\end{itemize}

\subsection*{E. Cross-Domain Translation}
\begin{itemize}[leftmargin=*]
    \item \textbf{Draw a Picture}: Visualize 24 equally-spaced points on the unit circle. Raising to the 6th power multiplies the angle by 6. Points that map to the real axis (angles $0, \pi, 2\pi, \ldots$) are the ones we count
    \item \textbf{Recast the Problem}: 
    \begin{itemize}
        \item Algebraic: $z^6$ real $\Leftrightarrow$ $z^6 = \overline{z^6} = \overline{z}^6$ $\Leftrightarrow$ $z^6/\overline{z}^6 = 1$ $\Leftrightarrow$ $(z/\overline{z})^6 = 1$
        \item Number theoretic: $z^6$ real $\Leftrightarrow$ argument of $z^6$ is $0 \pmod{\pi}$ $\Leftrightarrow$ $6 \cdot \frac{2\pi k}{24} \equiv 0 \pmod{\pi}$ $\Leftrightarrow$ $k$ even
        \item Group theoretic: We want $z^6 \in \langle -1 \rangle$, the subgroup of order 2 in the unit circle
    \end{itemize}
    \item \textbf{Change Your Point of View}: Instead of checking all 24 roots, recognize that $z^6$ being real means $z^6 \in \{1, -1\}$. So we're counting solutions to $z^6 = 1$ plus solutions to $z^6 = -1$ among the 24th roots of unity. Each equation has 6 solutions, but we must check they're all 24th roots of unity (they are)
\end{itemize}
\end{strategicbox}

\begin{tacticalbox}
\subsection*{A. Core Mathematical Tactics}
\begin{itemize}[leftmargin=*]
    \item \textbf{Symmetry Exploitation}: The 24th roots have rotational symmetry. If $z$ is a root with $z^6$ real, then $z^{-1}$ is also a root with $(z^{-1})^6 = (z^6)^{-1}$ real
    \item \textbf{Extreme Principle}: Not directly applicable
    \item \textbf{Invariant}: The property "$z^6$ is real" is preserved under the transformation $z \mapsto z \cdot e^{2\pi i/12}$ (jumping by 2 positions)
    \item \textbf{Monovariant}: Not applicable
    \item \textbf{Parity}: Key tactic! The index $k$ must be even for $z^6$ to be real. This is a parity condition
    \item \textbf{Coloring}: Can "color" the 24 roots alternately as even/odd indexed. Even indices (12 roots) give real $z^6$
    \item \textbf{Telescoping}: Not applicable
    \item \textbf{Generating Functions}: Not applicable to this discrete counting problem
\end{itemize}

\subsection*{B. Domain-Specific Tactics}
\begin{itemize}[leftmargin=*]
    \item \textbf{Complex Numbers}: 
    \begin{itemize}
        \item Exponential form: $z = e^{i\theta}$, $z^n = e^{in\theta}$
        \item Real condition: $\sin(\theta) = 0$ $\Leftrightarrow$ $\theta = n\pi$
        \item De Moivre's Theorem application
    \end{itemize}
    \item \textbf{Roots of Unity}: 
    \begin{itemize}
        \item The $n$-th roots of unity are $e^{2\pi i k/n}$ for $k = 0, 1, \ldots, n-1$
        \item They form a cyclic group under multiplication
        \item Subgroup structure: 24th roots contain 12th, 6th, 4th, 3rd, 2nd roots as subgroups
    \end{itemize}
    \item \textbf{Modular Arithmetic}: $k/2 \in \mathbb{Z}$ $\Leftrightarrow$ $k \equiv 0 \pmod{2}$
\end{itemize}

\subsection*{C. Cross-Domain Techniques}
\begin{itemize}[leftmargin=*]
    \item \textbf{Coordinate Geometry}: Unit circle representation
    \item \textbf{Trigonometry}: Using $e^{i\theta} = \cos\theta + i\sin\theta$
    \item \textbf{Abstract Algebra}: Group theory perspective (cyclic groups, subgroups, cosets)
\end{itemize}
\end{tacticalbox}

\begin{techniquesbox}
\subsection*{A. Recognition Index}
\begin{itemize}[leftmargin=*]
    \item \textbf{Signature Patterns Identified}:
    \begin{itemize}
        \item Keywords: "complex numbers", "$z^{24} = 1$", "roots of unity"
        \item Structure: Counting problem with algebraic constraint
        \item Answer choices: 12 is exactly half of 24 (strong parity hint)
    \end{itemize}
    \item \textbf{Instant Recognition Triggers}: "$z^n = 1$" immediately suggests $n$-th roots of unity
    \item \textbf{Quick Classification}: Late-band AMC12 complex number problem, standard difficulty for students familiar with roots of unity
\end{itemize}

\subsection*{B. Technique Selection Process}
\begin{itemize}[leftmargin=*]
    \item \textbf{Primary Technique}: Exponential form of complex numbers with parity analysis
    \begin{itemize}
        \item Probability of success: 90-95\% (standard technique for this problem type)
        \item Time estimate: 2-3 minutes
        \item Required knowledge: Euler's formula $e^{i\theta} = \cos\theta + i\sin\theta$, roots of unity
    \end{itemize}
    \item \textbf{Alternative Techniques}:
    \begin{itemize}
        \item Trigonometric form: $z = \cos(2\pi k/24) + i\sin(2\pi k/24)$, then compute $z^6$ using De Moivre
        \item Algebraic substitution: $z^6$ real means $z^6 \in \{1, -1\}$, count 6th roots of 1 and 6th roots of $-1$
        \item Answer choice elimination: test specific values and identify pattern
    \end{itemize}
    \item \textbf{Technique Confidence Ranking}:
    \begin{enumerate}
        \item Exponential form + parity (highest confidence, most direct)
        \item Algebraic: $z^6 \in \{1,-1\}$ counting approach (equally valid)
        \item Trigonometric computation (more computational, same answer)
        \item Geometric visualization (good for intuition, less rigorous)
    \end{enumerate}
\end{itemize}

\subsection*{C. Integration Strategy}
\begin{itemize}[leftmargin=*]
    \item \textbf{Multi-Technique Combination}: 
    \begin{itemize}
        \item Use exponential form for primary computation
        \item Apply parity argument for counting
        \item Verify with alternative $z^6 \in \{1, -1\}$ approach
    \end{itemize}
    \item \textbf{Technique Transition}: If exponential form is unfamiliar, fall back to testing small cases and identifying the pattern
    \item \textbf{Verification Bridge}: The two approaches (parity and $z^6 \in \{1,-1\}$) serve as mutual verification
\end{itemize}

\subsection*{D. Conflict Resolution}
\begin{itemize}[leftmargin=*]
    \item \textbf{Potential Conflicts}: 
    \begin{itemize}
        \item Confusion between "real" and "positive real" - ensure understanding that negative reals count
        \item Off-by-one errors in counting even numbers from 0 to 23
        \item Forgetting that $k=0$ is even (giving the root $z=1$)
    \end{itemize}
    \item \textbf{Resolution Strategy}: Explicitly list the even values: $0, 2, 4, 6, 8, 10, 12, 14, 16, 18, 20, 22$ (12 values)
    \item \textbf{Compatibility Check}: Both approaches (parity and algebraic) must yield 12 as the answer
\end{itemize}
\end{techniquesbox}

\begin{verificationbox}
\subsection*{A. Verification Level Selection}

\textbf{Recommended Level}: Level 4 (Standard) - This is a late-band problem worth verifying carefully

\textbf{Decision Tree}:
\begin{itemize}[leftmargin=*]
    \item Confidence: 85-90\% (standard technique, straightforward execution)
    \item Problem position: P17 (late-band, important to get right)
    \item Time available: 2-3 minutes for verification appropriate
    \item Consequence of error: Missing a late-band problem significantly impacts score
\end{itemize}

\subsection*{B. Specialized Verification Methods}

\textbf{Complex Numbers Verification}:
\begin{itemize}[leftmargin=*]
    \item \textbf{Direct computation verification}: Test several even and odd values of $k$
    \begin{itemize}
        \item $k=0$ (even): $z^6 = e^0 = 1$ (real) \checkmark
        \item $k=1$ (odd): $z^6 = e^{i\pi/2} = i$ (not real) \checkmark
        \item $k=2$ (even): $z^6 = e^{i\pi} = -1$ (real) \checkmark
        \item $k=3$ (odd): $z^6 = e^{3i\pi/2} = -i$ (not real) \checkmark
    \end{itemize}
    \item \textbf{Algebraic verification}: Count solutions to $z^6 = 1$ and $z^6 = -1$ among 24th roots
    \begin{itemize}
        \item $z^6 = 1$: Need $e^{i\pi k/2} = 1$, so $k \in \{0, 4, 8, 12, 16, 20\}$ (6 values)
        \item $z^6 = -1$: Need $e^{i\pi k/2} = -1$, so $k \in \{2, 6, 10, 14, 18, 22\}$ (6 values)
        \item Total: $6 + 6 = 12$ \checkmark
    \end{itemize}
    \item \textbf{Geometric verification}: On the unit circle, 24 points spaced $15^\circ$ apart. After rotating by $6 \times 15^\circ = 90^\circ$, which land on real axis? Every other point alternately
\end{itemize}

\textbf{Answer Choice Verification}:
\begin{itemize}[leftmargin=*]
    \item Answer (D) 12 makes sense: exactly half the roots, consistent with parity argument
    \item Answer (B) 4 or (C) 6 would require different divisibility pattern (ruled out)
    \item Answer (E) 24 would mean all $z^6$ are real (clearly false, e.g., $z = e^{i\pi/12}$ gives $z^6 = i$)
    \item Answer (A) 0 would mean no $z^6$ is real (false, e.g., $z=1$ gives $z^6=1$)
\end{itemize}

\subsection*{C. Confidence Calibration}

\textbf{Confidence Assessment}: 95\%

\textbf{Justification}:
\begin{itemize}[leftmargin=*]
    \item Two independent methods both yield 12
    \item Parity argument is clean and rigorous
    \item Algebraic verification ($z^6 \in \{1, -1\}$) confirms the count
    \item Answer choice 12 is exactly half, matching our even/odd pattern
    \item No computational errors in exponential form calculations
    \item Small case testing confirms the pattern
\end{itemize}

\textbf{Remaining Uncertainty}: 5\% from potential misunderstanding of "real number" definition or counting error
\end{verificationbox}

\begin{solutionbox}
\subsection*{Complete Solution Roadmap}

\textbf{Step 1: Identify the 24th roots of unity} (30 seconds)

The solutions to $z^{24} = 1$ are the 24th roots of unity:
\[z_k = e^{2\pi i k/24} = e^{i\pi k/12}\]
for $k = 0, 1, 2, \ldots, 23$.

\textbf{Step 2: Compute $z^6$ for a general root} (30 seconds)

For the $k$-th root $z_k = e^{i\pi k/12}$, we have:
\[z_k^6 = \left(e^{i\pi k/12}\right)^6 = e^{i\pi k/2}\]

\textbf{Step 3: Determine when $z^6$ is real} (1 minute)

A complex number $e^{i\theta}$ is real if and only if $\sin(\theta) = 0$, which occurs when $\theta$ is an integer multiple of $\pi$.

For $z_k^6 = e^{i\pi k/2}$ to be real, we need:
\[\frac{\pi k}{2} = n\pi \text{ for some integer } n\]

This simplifies to:
\[\frac{k}{2} = n \quad \Rightarrow \quad k = 2n\]

Therefore, $z_k^6$ is real if and only if $k$ is even.

\textbf{Step 4: Count the even values} (30 seconds)

The even values of $k$ in the range $\{0, 1, 2, \ldots, 23\}$ are:
\[k \in \{0, 2, 4, 6, 8, 10, 12, 14, 16, 18, 20, 22\}\]

There are exactly $\boxed{12}$ such values.

\textbf{Alternative Verification}: We can verify this by noting that $z^6$ real means $z^6 \in \{1, -1\}$ (the only real numbers on the unit circle). 
\begin{itemize}
    \item Solutions to $z^6 = 1$: $k \in \{0, 4, 8, 12, 16, 20\}$ (6 values)
    \item Solutions to $z^6 = -1$: $k \in \{2, 6, 10, 14, 18, 22\}$ (6 values)
    \item Total: $6 + 6 = 12$ \checkmark
\end{itemize}

\subsection*{Checkpoints and Alternatives}

\textbf{Checkpoint 1} (After Step 2): Verify the exponent calculation
\begin{itemize}[leftmargin=*]
    \item Confirm: $(e^{i\pi k/12})^6 = e^{i \cdot 6\pi k/12} = e^{i\pi k/2}$
    \item If error here, revisit exponent rules
\end{itemize}

\textbf{Checkpoint 2} (After Step 3): Verify the real condition
\begin{itemize}[leftmargin=*]
    \item Test case: $k=2$ gives $e^{i\pi} = -1$ (real) \checkmark
    \item Test case: $k=1$ gives $e^{i\pi/2} = i$ (not real) \checkmark
\end{itemize}

\textbf{Alternative Approach}: Geometric visualization
\begin{itemize}[leftmargin=*]
    \item The 24 roots form a regular 24-gon on the unit circle
    \item Raising to 6th power rotates each point by $6 \times (360^\circ/24) = 90^\circ$
    \item Points initially at $0^\circ, 30^\circ, 60^\circ, \ldots$ map to $0^\circ, 180^\circ, 360^\circ, \ldots$
    \item Only points at multiples of $180^\circ$ after rotation are real
    \item This occurs for the 12 even-indexed original points
\end{itemize}

\subsection*{Comprehensive Pitfall Guide}

\begin{warningbox}
\textbf{Pitfall 1}: Confusing "real" with "positive real"

Negative real numbers like $-1$ are still real! The problem asks for real, not positive real.

\textbf{Prevention}: Remember that $e^{i\pi} = -1$ is real.
\end{warningbox}

\begin{warningbox}
\textbf{Pitfall 2}: Off-by-one error in counting

Forgetting that $k=0$ is even, or that $k$ ranges from 0 to 23 (not 1 to 24).

\textbf{Prevention}: Explicitly list: $0, 2, 4, 6, 8, 10, 12, 14, 16, 18, 20, 22$ and count.
\end{warningbox}

\begin{warningbox}
\textbf{Pitfall 3}: Incorrect application of De Moivre's Theorem

Writing $z^6 = e^{i\pi k/12}$ instead of $z^6 = e^{i \cdot 6 \cdot \pi k/12} = e^{i\pi k/2}$.

\textbf{Prevention}: Carefully multiply the exponent: $(e^{i\alpha})^n = e^{in\alpha}$.
\end{warningbox}

\begin{warningbox}
\textbf{Pitfall 4}: Misidentifying when a complex number is real

Thinking that $e^{i\theta}$ is real when $\theta$ is rational, or when $\cos\theta \neq 0$.

\textbf{Prevention}: $e^{i\theta} = \cos\theta + i\sin\theta$ is real $\Leftrightarrow$ $\sin\theta = 0$ $\Leftrightarrow$ $\theta = n\pi$.
\end{warningbox}

\subsection*{Time Allocation Breakdown}
\begin{itemize}[leftmargin=*]
    \item Problem reading and setup: 30 seconds
    \item Identify roots of unity form: 30 seconds
    \item Compute $z^6$ expression: 30 seconds
    \item Determine real condition: 1 minute
    \item Count even values: 30 seconds
    \item Verification: 1 minute
    \item \textbf{Total: 4 minutes}
\end{itemize}
\end{solutionbox}

\begin{competitionbox}
\subsection*{A. Time Management Strategy}

\textbf{Time Budget}: 3-4 minutes for P17 (late-band)

\textbf{Actual Time}: 4 minutes (primary solution) + 1 minute (verification) = 5 minutes total

\textbf{Assessment}: Appropriate time investment for a late-band problem. If unfamiliar with roots of unity, this problem should be skipped initially and returned to if time permits.

\subsection*{B. Problem Prioritization}

\textbf{Priority Level}: Medium-High
\begin{itemize}[leftmargin=*]
    \item P17 is challenging but accessible with right knowledge
    \item If comfortable with complex numbers: solve immediately
    \item If unfamiliar: skip and solve easier P18-20 first, return if time permits
\end{itemize}

\textbf{Domain Comfort Check}:
\begin{itemize}[leftmargin=*]
    \item Do you know Euler's formula $e^{i\theta} = \cos\theta + i\sin\theta$? 
    \item Do you recognize "$z^{24}=1$" as roots of unity?
    \item If yes to both: proceed. If no: consider skipping.
\end{itemize}

\subsection*{C. Risk Management}

\textbf{Confidence Threshold}: 95\%

\textbf{Decision}: Submit answer (D) 12 with high confidence and move on.

\textbf{Flag for Review}: No - verification passed, two independent methods agree, answer choice makes sense.

\textbf{Error Recovery Protocol}: If discovered error in verification (e.g., miscounted even numbers):
\begin{itemize}[leftmargin=*]
    \item Recount: $0, 2, 4, 6, 8, 10, 12, 14, 16, 18, 20, 22$ = 12 values
    \item Alternative count: There are $\lfloor 23/2 \rfloor + 1 = 12$ even numbers from 0 to 23
\end{itemize}

\subsection*{D. Abandonment Criteria}

\textbf{Soft Abandon Triggers} (consider moving on):
\begin{itemize}[leftmargin=*]
    \item After 3 minutes, if roots of unity concept is completely unfamiliar
    \item If unable to express $z^6$ in exponential form after 2 minutes
    \item If answer doesn't match any answer choice after careful calculation
\end{itemize}

\textbf{Hard Abandon Triggers} (definitely switch problems):
\begin{itemize}[leftmargin=*]
    \item Realize you don't know what complex numbers are
    \item Getting emotionally frustrated with exponential form after 4 minutes
    \item Made same calculation error twice
\end{itemize}

\textbf{Strategic Abandonment}: 
\begin{itemize}[leftmargin=*]
    \item If abandoning, make educated guess: answer (D) 12 is most likely given it's exactly half of 24
    \item Mark problem to return if time permits
    \item Move to P18 which may be more accessible
\end{itemize}
\end{competitionbox}

\begin{keyinsightbox}
\textbf{Key Insight}: The 24th roots of unity can be written as $z = e^{i\pi k/12}$ for $k = 0, 1, \ldots, 23$. Then $z^6 = e^{i\pi k/2}$, which is real precisely when $\pi k/2$ is an integer multiple of $\pi$, i.e., when $k$ is even. There are exactly 12 even values of $k$ in the range, so the answer is 12.

\textbf{Pattern Recognition}: This problem tests whether you recognize the parity pattern in roots of unity. The answer being exactly half the total number of roots (12 out of 24) is a strong hint that an even/odd argument is involved.

\textbf{Alternative Perspective}: Since $z^6$ being real means $z^6 \in \{1, -1\}$, we're essentially counting how many 24th roots of unity are also 6th roots of 1 or 6th roots of $-1$. Each set has 6 elements, giving $6 + 6 = 12$ total.
\end{keyinsightbox}

\begin{tipbox}
\textbf{Pro Tip for Complex Numbers}: Whenever you see "$z^n = 1$", immediately think "$n$-th roots of unity" and write them in exponential form $e^{2\pi i k/n}$. This standard form makes all calculations straightforward.

\textbf{Competition Strategy}: For late-band problems (P17-25), if you don't see the solution path within 1-2 minutes, check if the answer choices give hints. Here, 12 being exactly half of 24 strongly suggests a parity or symmetry argument.

\textbf{Verification Tip}: Always test both an even and odd case when you suspect a parity pattern. This quickly confirms or refutes your hypothesis.
\end{tipbox}

\section*{Final Answer}
\boxed{\textbf{(D)}\ 12}

\section*{Confidence Assessment}
\begin{itemize}[leftmargin=*]
    \item \textbf{Confidence Level}: 95\% - Standard technique properly applied, verified with alternative method
    \item \textbf{Verification Applied}: Level 4 (Standard) - Tested specific cases, used alternative algebraic approach, answer choice analysis
    \item \textbf{Flag for Review}: No - High confidence, multiple verification methods agree
    \item \textbf{Time Spent}: 5 minutes (4 min solution + 1 min verification) - Appropriate for late-band problem
\end{itemize}

\end{document}